\documentclass[10pt]{article}
\usepackage[a4paper, margin=1.25cm]{geometry}
\usepackage{
  fontspec,
  multicol,
  graphicx,
  enumitem,  % for \begin{itemize}[noitemsep, nolistsep] etc.
  hanging,
  tabularx,
  dblfnote   % also contained in package 'yafoot' (install that if need be)
}
\usepackage[hidelinks]{hyperref}

\setmainfont{Avenir}
\setmonofont{Courier}

\setlength{\parskip}{0.33em}

\pagestyle{empty}
\raggedright

\DFNalwaysdouble  % 2-columns of footnotes
\renewcommand{\footnoterule}{% full-width footnote rule 
  \kern -3pt
  \hrulefill \break
  \kern 2pt
}
\newcommand{\myfootnotesize}{\footnotesize}
\renewcommand{\footnotesize}{\scriptsize}



\begin{document}

\begin{multicols}{3}
  \
  \vfill
  \raggedleft
    {\huge \textbf{Ruan van Mazijk}} \par
    {\small Ecologist, Editorial \& Training Specialist} \par
    C4 EcoSolutions (Pty) Ltd.
  \vfill
  \
  \begin{center}
    \includegraphics[width=8em]{IMG_4958_crop.png}
  \end{center}
  \columnbreak
  \par \
  \vfill
  \raggedright
    {\small \href{mailto:ruanvmazijk@gmail.com}{ruanvmazijk@gmail.com}} \par
    {\small (+27) 72 516 7570} \par
    {\small @rvanmazijk}
      \href{https://github.com/rvanmazijk}
           {\includegraphics[width=1em]{logos/GitHub.png}}
      \href{https://orcid.org/0000-0003-2659-6909}
           {\includegraphics[width=1em]{logos/ORCID.png}}
      \href{https://twitter.com/rvanmazijk}
           {\includegraphics[width=1em]{logos/Twitter.png}}
      \href{https://www.researchgate.net/profile/Ruan-Van-Mazijk}
           {\includegraphics[width=1em]{logos/ResearchGate-white_crop.jpg}}
  \vfill
  \
\end{multicols}

\hrule \vskip2pt \hrule
 % =======================================================

\section*{Education} % ---------------------------------------------------------

\begin{hangparas}{2em}{1}

\textbf{MSc candidate} in Biological Sciences
  {\small (UCT)                                                               \\
  \hspace{2em} \textit{Genome-size effects on plant ecophysiology \&          \\
  \hspace{4em} phenology in Cape schoenoid sedges}                            \\
  \hspace{2em} Principal supervisor: Prof A.M.~Muasya                         \\
  \hspace{2em} Co-supervisors: Prof G.A.~Verboom \& Prof A.G.~West}

\textbf{BSc Hons} in Biological Sciences {\small (UCT) \textit{With distinction} \hfill 2017 \\
  \hspace{2em} Supervisors: Prof M.D.~Cramer \& Prof G.A.~Verboom}

\textbf{BSc} in Ecology \& Evolution, Applied Biology {\small (UCT) \hfill 2016 \\
  \hspace{2em} \textit{With distinction in both majors \& the degree}}

\end{hangparas}


\section*{Publications} % ------------------------------------------------------

Ismail, Duker, \textbf{van Mazijk}, Mills \& Potts.
The importance of ground-truthing for thicket restoration    \hfill In prep. \\
\hspace{2em} site-selection: distinguishing degraded solid thicket from natural thicket mosaics. \textit{Ecol.~Rest.}

Ismail, \textbf{van Mazijk}, Duker \& Mills.
\textit{Portulacaria~afra} canopies support diverse understory plant growth-forms
                                                              \hfill In prep. \\
\hspace{2em} in southern African xeric succulent thicket. \textit{J.~Arid~Env.}

Mills \& \textbf{van Mazijk}.
Permutation tests to identify significant constraint or promotion within biological \hfill 2024 \\
\hspace{2em} scatterplots.
  \textit{Ecol.~Evol.}
  14. \href{http://dx.doi.org/10.1002/ece3.70584}{10.1002/ece3.70584}

Mills, \textbf{van Mazijk}, Allen \& Strydom.
Soil geochemistry \& constraint of tree seedlings immediately after
                                                              \hfill In rev. \\
\hspace{2em} germination on \textit{Macrotermes} termite
  mounds in the Kruger National Park, South Africa.
  \textit{Oikos}

\textbf{Van Mazijk}, West, Verboom, Elliott, Bureš \& Muasya.
Genome size variation in Cape schoenoid sedges                \hfill In rev. \\
\hspace{2em} (Cyperaceae: Schoeneae) \& its ecophysiological consequences.
  \textit{Am.~J.~Bot.}

Mills, Duker, Lechmere-Oertel, \textbf{van Mazijk} \& Potts.
Science is the fuel required for lifting ecosystem             \hfill In press \\
\hspace{2em} restoration into the orbit of hundreds of millions of hectares.
  \textit{Front.~Cons.~Sci.}
  4. \\
  \hspace{2em} \href{https://doi.org/10.3389/fcosc.2023.1283262}{10.3389/fcosc.2023.1283262}

\textbf{Van Mazijk}, Cramer \& Verboom.
Environmental heterogeneity explains contrasting plant species richness
                                                                  \hfill 2021 \\
\hspace{2em}
  between the South African Cape \& southwestern Australia.
  \textit{J.~Biogeog.}
  48. \href{https://doi.org/10.1111/jbi.14118}{10.1111/jbi.14118}.

Elliott, \textbf{van Mazijk}, Barrett, Bruhl,
Joly, Muthaphuli, Wilson \& Muasya.
Global dispersal \& diversification of                            \hfill 2021 \\
\hspace{2em}
  the genus \textit{Schoenus} (Cyperaceae)
  from the Western Australian biodiversity hotspot.                           \\
\hspace{2em}
  \textit{J.~Syst.~Evol.}
  59(4).
  \href{https://doi.org/10.1111/jse.1274}{10.1111/jse.12742}.

Elliott, Laidler, Muasya, Muthaphuli \& \textbf{van Mazijk}.
Solving the sedge mystery.
\textit{Veld \& Flora}~105(3):20--25.                             \hfill 2019

Cramer, Wootton, \textbf{van Mazijk} \& Verboom.
New regionally-modelled soil layers improve prediction of         \hfill 2019 \\
\hspace{2em}
  vegetation type relative to that based on global soil models.
  \textit{Div.~Distrib.}
  25(11).                                                        
  \href{https://doi.org/10.1111/ddi.12973}{10.1111/ddi.12973}.

\textbf{Van Mazijk}, Smyth, Weideman \& West.
Isotopic tracing of stormwater in the urban Liesbeek River.
\textit{Water~SA}                                                 \hfill 2018 \\
\hspace{2em}
  44(4).
  \href{https://doi.org/10.4314/wsa.v44i4.16}{10.4314/wsa.v44i4.16}.


\section*{Peer review} % -------------------------------------------------------

Grobler \& Cowling. The composition, geography, biology and assembly of the
  coastal flora of the Cape                                       \hfill 2021 \\
  \hspace{2em} Floristic Region. \textit{PeerJ} 9:e11916. DOI:
    \href{https://doi.org/10.7717/peerj.11916}{10.7717/peerj.11916}.

\hfill

\begin{multicols}{2}

\section*{Interests} % ---------------------------------------------------------
    
Systematics \& taxonomy

Phylogenetic comparative methods

Biodiversity patterns \& macroecology

Community ecology

Plant ecophysiology \& functional traits

\section*{Professional affiliations} % -----------------------------------------

SASSB\footnote{SASSB: Southern African Society for Systematic Biology}
council student representative                                 \hfill 2019--2021

SASSB student member                                               \hfill 2018--

SAAB\footnote{SAAB: South African Association of Botanists}
student member                                                     \hfill 2017--

Golden Key International Honour Society                            \hfill 2015--

\end{multicols}

\vfill

\clearpage % ===================================================================

\section*{Funding \& awards} % -------------------------------------------------

\subsection*{Travel grants}

SASSB \& GENUS (DSI-NRF Centre of Excellence in Palaeosciences) Sponsorship for SASSB XVI (R6,000)                                   \hfill 2023

NRF Student Sponsorship for online attendance at MEDECOS XVII
  (R7,000)                                                           \hfill 2022

UCT Travel Grant for the 45\textsuperscript{th} Annual SAAB,
  AMA\footnote{AMA: African Mycological Assocation} \& SASSB Joint
  Congress, 2019 (R3,819)                                            \hfill 2018

\subsection*{Scholarships \& bursaries}

UCT Postgraduate Publication Incentive Funding (R4,000)              \hfill 2023

NRF Innovation Masters Scholarship (R90,000)                         \hfill 2019

UCT Masters Research Scholarship\footnote{Awarded but not reimbursed
  due to award of Innovation Masters Scholarship} (R30,000)          \hfill 2019

UCT Bursary Top Ups\footnote{From MSc supervisors} (R6,000 total)    \hfill 2019

SAAB MSc Student Bursary (inaugural recipient) (R20,000 p.a.)  \hfill 2018--2019

UCT Department of Biological Sciences Dorothy Cameron Scholarship
  (R20,000)                                                          \hfill 2018

UCT Once-Off Top Up Award (R5,520)                                   \hfill 2018

UCT Masters Research Scholarship (R30,000)                           \hfill 2018

NRF Innovation Honours Scholarship (R60,000)                         \hfill 2017

SAAB Honours Scholarship (R10,000)                                   \hfill 2017

UCT Council Honours Merit Scholarship\footnote{Awarded but not
  reimbursed due to award of Innovation Honours Scholarship}
  (R20,000)                                                          \hfill 2017

UCT Science Faculty Scholarship (R15,000 p.a.)                 \hfill 2015--2016

\subsection*{Awards \& honours}

3\textsuperscript{rd} Best Student Presentation at SASSB XVI     \hfill Jul 2023

Best Systematics MSc Presentation\footnote{Awarded by SASSB,
  3-way tie with Z.~Shaik (UCT, SANBI) \& D.~Main (UJ)}
  at the 45\textsuperscript{th}
  Annual SAAB, AMA \& SASSB Joint Congress                       \hfill Jan 2019

UCT Dean's Merit List                                          \hfill 2014--2016

UCT Dick \& Dorothy Borcherds Prize, for the highest standard in 
  2\textsuperscript{nd} year biology (R1,000)                        \hfill 2015

Top of the class at UCT for:

\hspace{2em} 3\textsuperscript{rd} year evolutionary biology         \hfill 2016

\hspace{2em} 2\textsuperscript{nd} year botany                       \hfill 2015 

\hspace{2em} 2\textsuperscript{nd} year zoology                      \hfill 2015

\hspace{2em} 2\textsuperscript{nd} year marine biology               \hfill 2015

\hspace{2em} 1\textsuperscript{st} year biological diversity         \hfill 2014

\section*{Leadership positions} % ----------------------------------------------

SASSB council member, student representative                   \hfill 2019--2021

UCT Department of Biological Sciences Postgraduate Committee member
                                                           \hfill Oct 2018--2019

UCT Departmental Representative for the Department of Biological Sciences
                                                                     \hfill 2016

Class Representative at UCT for:

\hspace{2em} 3\textsuperscript{rd} year evolutionary biology         \hfill 2016

\hspace{2em} 3\textsuperscript{rd} year ecology                      \hfill 2016

\hspace{2em} 2\textsuperscript{nd} year botany                       \hfill 2015

\clearpage % ===================================================================

\section*{Teaching assistance} % -----------------------------------------------

% TODO: tidy this up to be more **pow** like my 1p CV

\begin{hangparas}{2em}{1}

Teaching assistant \& field assistant for 2\textsuperscript{nd} year angiosperm
diversity                                                     \hfill 2022 \break
Convenor: Associate Prof S.~Chimphango (UCT)

Teaching \& field assistant for 2\textsuperscript{nd} year ecology \& evolution
                                                        \hfill 2019--2020 \break
Convenor: Prof G.A.~Verboom (UCT)

Teaching assistant for 2\textsuperscript{nd} year angiosperm diversity
                                                        \hfill 2018--2019 \break
Lecturer: Prof G.A.~Verboom (UCT)

Tutor for 2\textsuperscript{nd} year study design \& data analysis for 
scientists in \textit{R}                                \hfill 2017--2020 \break
Convenor: Associate Prof R.~Altwegg (SEEC\footnote{SEEC: Centre for 
Statistics in Ecology, the Environment \& Conservation, UCT})
% NOTE: finished tutoring for STA2007* in Aug 2020

Lectured for SDU\footnote{SDU: Schools Development Unit, UCT} short course 
\textit{Teaching Biodiversity for FET\footnote{Further Education \& Training
(grades 10--12 in South Africa)}} Life Sciences Teachers
                                                       \hfill 18 Mar 2019 \break
Lecture title: \textit{Biodiversity, evolution \& taxonomy---an overview of the 
major clades of animals \& plants}

Assistant tutor for SEEC-ACCESS\footnote{ACCESS: Applied Centre for Climate \&
Earth Systems, NRF} short course \textit{Study design \& data analysis for 
scientists in} R                                   \hfill 18--20 Jul 2018 \break
Organiser: Associate Prof R.~Altwegg (SEEC)

Assistant tutor for 1\textsuperscript{st} year introductory statistics
\textit{R} practicals                                   \hfill 2016--2018 \break
Convenor: Associate Prof R.~Altwegg (SEEC)

Hot-seat tutor for 1\textsuperscript{st} year mathematics for biology, earth \& 
environmental sciences undergraduates                         \hfill 2015 \break
Convenor: T.~van Heerden (UCT); hot-seat organiser: Dr R.~Moolman (UCT)

\hfill

\end{hangparas}

\section*{Workshops} % ---------------------------------------------------------

\begin{hangparas}{2em}{1}

Organiser, lead for \textit{Data wrangling \& manipulation in R}
                                                     \hfill 1--3 Jul 2019 \break
For UCT Department of Biological Sciences postgraduates
% TODO: add URL

Attendee at \textit{Phylogenetics in R}              \hfill 6--8 Jan 2019 \break
Pre-congress workshop at the 45th Annual SAAB, AMA \& SASSB Joint Congress\break
Instructors:
  Prof L.J.~Revell (University of Massachusetts Boston)                   \break
  \& Prof L.J.~Harmon (University of Idaho)

Presenter at \textit{Biodiversity \& Environmental Change in the Cape 
Floristic Region}                                      \hfill 24 Jan 2017 \break
With L.M. Wootton                                                         \break
Progress report on \textit{Development of soil atlas for the 
Greater Cape Floristic Region}                                            \break
Organiser: Associate Prof L.~Gillson (UCT)

Attendee at \textit{Spatial Point Process Modelling for Ecological 
Survey Data with} inlabru                          \hfill 11--13 Jan 2017 \break
Instructors:
  Prof D.L.~Borchers (CREEM\footnote{CREEM: Centre for Research into Ecological \& Environmental Modelling, St Andrews University}),
  Dr J.B.~Illian (University of Glasgow)                                \break
  \& Dr F.~Bachl (University of Edinburgh)

\hfill

\end{hangparas}

\section*{Research assistance} % -----------------------------------------------

\begin{hangparas}{2em}{1}

Field assistant for sampling \textit{Schoenus} spp.~(Cyperaceae, tribe 
Schoeneae)                                                    \hfill 2018 \break
Specimens were used for genome-size analyses and taxonomic study          \break
Supervisor: Dr T.L.~Elliott (Masaryk University) \& Prof A.M.~Muasya (UCT)

Laboratory assistant for \textit{Schoenus} spp.~(Cyperaceae, tribe Schoeneae) 
hydroponics pilot study                                       \hfill 2018 \break
Supervisor: Dr T.L.~Elliott (Masaryk University) \& Prof A.M.~Muasya (UCT)

Data processing assistant for assessing the current available data for the 
soils of the GCFR                                    \hfill Jun--Dec 2016 \break
With L.M.~Wootton                                                         \break
\textit{Biodiversity and Environmental Change in the Cape Floristic Region} 
research programme                                                        \break
Supervisor: Prof M.D.~Cramer (UCT)

Field assitant for population monitoring of
  \textit{Mimetes} spp.~\&
  \textit{Leucospermum} spp.~(Proteaceae)            \hfill Jan--Feb 2017 \break
  on the Cape Peninsula                                                   \break
With L.M.~Wootton                                                         \break
Supervisor: Dr J.A.~Slingsby (UCT, SAEON)

Laboratory \& field assistant for legume symbiotic rhizobia culturing
                                                          \hfill Dec 2015 \break
Also involved plant collections \& literature-derived rhizobial collection 
data processing                                                           \break
Supervisors: Dr M.~Dludlu (University of Swaziland) \& Prof A.M.~Muasya (UCT)

% I contributed to the West Lab/ACCESS/SEEC data management policy document

\hfill

\end{hangparas}

\clearpage % ===================================================================

\section*{Conferences} % -------------------------------------------------------

\begin{hangparas}{2em}{1}

Presenter at SASSB XVI                               \hfill 5--7 Jul 2023 \break
\textit{Genome size variation in Cape schoenoid sedges (Cyperaceae: Schoeneae) \& \break
its ecophysiological consequences}

Presenter at UCT Biological Sciences Research Day      \hfill 29 Nov 2022 \break
\textit{Genome size \& ecophysiology in two species-rich Cape lineages}

Presenter at SAAB Postgraduate Virtual Symposium       \hfill 28 Sep 2022 \break
\textit{Comparing genome size \& ecophysiology in two species-rich Cape sedge lineages \break
(}Schoenus\textit{ L.~\& }Tetraria\textit{ P.Beauv., tribe Schoeneae, Cyperaceae)}

Attendee at MEDECOS XVII                                    \hfill 5--9 Sep 2022

Presenter at IBS\footnote{IBS: International Biogeography Society}
Early Career Biogeographers Conference                    \hfill 22--24 Oct 2021 \\
\textit{Environmental heterogeneity explains contrasting plant species richness between \break
the South African Cape \& southwestern Australia}

Presenter \& candidate at iDiv\footnote{iDiv: German Centre for Integrative
Biodiversity Research, Halle-Jena-Leipzig} Online Joint Recruitment Symposium
                                                       \hfill 31 Aug--1 Sep 2020

Presenter at the 45th Annual SAAB, AMA \& SASSB Joint Congress
                                                        \hfill 9--11 Jan 2019 \\
\hspace{2em} 1\textsuperscript{st} oral presentation:       \hfill 9 Jan 2019 \\
\hspace{4em} \textit{Environmental turnover predicts plant species richness \& 
turnover}                                                                     \\
\hspace{4em} Based on my then BSc Hons work, later published in \textit{Journal
of Biogeography}                                                              \\
\hspace{2em} 2\textsuperscript{nd} oral presentation:      \hfill 11 Jan 2019 \\
\hspace{4em} \textit{Larger genome size constrains water-use efficiency \&
phenology in Cape schoenoid sedges}                                           \\
\hspace{4em} Based on my MSc work

Presenter at \textit{UCT BioDay}                       \hfill 30 Nov 2018 \break
Preliminary results of my MSc work

Attendee at \href{https://capetown2018.satrdays.org}{\textit{satRday Cape Town}}
                                                              \hfill 17 Mar 2018

\hfill

\end{hangparas}

\section*{Press releases \& science communication} % ---------------------------

\begin{hangparas}{2em}{1}

% TODO: add podcast invite

% TODO: add JBI ECR feature

Radio interview on \textit{The John Maytham Show}, \textit{CapeTalk}
                                                       \hfill 28 Nov 2018 \break
\href{https://www.capetalk.co.za/articles/328900/harvesting-stormwater-from-liesbeek-river-may-aid-ct-water-supply-students-find}
{\textit{Harvesting Storm Water}}

Interviewed for article for UCT News by S.~Segar       \hfill 26 Nov 2018 \break
\textit{Harvesting storm water from the Liesbeek River}
% TODO: add URL

Mentioned in article for UCT Faculty of Science press releases
                                                     \hfill Nov--Dec 2018 \break
\textit{Tracing stormwater in the Liesbeeck River - students passion drives
scientific discovery}
% TODO: add URL


Presenter at \href{http://pintofsciencesa.wixsite.com/pintofsciencesa/uct-planet-earth-16}
{\textit{Pint of Science Cape Town}}                   \hfill 16 May 2018 \break
I talked about hypotheses concerning the role of genome size in plant water use 
\& ecology                                                                \break
to a general audience

\hfill

\end{hangparas}

\begin{multicols}{2}

\section*{Skills} % ------------------------------------------------------------

Strong public speaker \\
Written communication skills \\
Teaching experience \\
Visual design \\
Detail oriented \\
Analytical thinking \\
Highly numerate

\section*{Methods} % -----------------------------------------------------------

Data manipulation, wrangling \& cleaning \\
Data visualisation \\
Statistical modelling \\
Phylogenetic inference \\
Phylogenetic comparative methods \\
GIS (in \textit{R})

\columnbreak

\section*{Languages} % ---------------------------------------------------------

English (1\textsuperscript{st} language) \\
Afrikaans (2\textsuperscript{nd} language)

\section*{Tools} % -------------------------------------------------------------

\textit{bash}, \textit{R} \\
\textit{R Markdown} \& related \textit{R}-packages \\
\textit{tidyverse} \textit{R}-packages \\
\textit{git}, GitHub \\
\textit{LaTeX}, \textit{BibTeX} \\
Apple macOS \\
Microsoft Windows \\
Microsoft Offce Word, Excel \& Powerpoint

\end{multicols}

\clearpage % ===================================================================

\section*{Miscellaneous} % -----------------------------------------------------

\begin{hangparas}{2em}{1}

UCT Open Day volunteer \hfill 22 Apr 2017, 21 Apr 2018, 4 May 2019, 8 Oct 2022 \break
For the Department of Biological Sciences

Animal Ethics Course certified by South African Veterinary Council
                                                          \hfill Apr 2017 \break
Completed as part of BSc Hons general coursework

Botanical Society of South Africa (BotSoc) member                  \hfill 2017--

UCT Wine \& Culture Society member                             \hfill 2016--2018

UCT Birding Club member                                        \hfill 2016--2017

Selected to give student feedback for curriculum-restructuring
                                                          \hfill Aug 2016 \break
Within the Department of Biological Sciences                              \break
Related to roles as Class/Departmental Representative---see above         \break
Selected by: Prof A.G.~West (UCT)

Selected to give student feedback for \textit{Biology: How Life Works}
(2\textsuperscript{nd} edition)                           \hfill May 2016 \break
Morris et al., 2016. W.H.~Freeman \& Co., Macmillan Education             \break
Selected by: Dr R.~Kelly-Laubscher (University College Cork)

\textit{Cape Wine Ambassador} qualification                   \hfill 2016 \break
Conferred by: Clive Torr

UCT Biological Society member                                        \hfill 2015

Two Oceans Aquarium \textit{Young Biologist} Volunteer         \hfill 2013--2014

\hfill

\end{hangparas}

\section*{References} % --------------------------------------------------------

\begin{hangparas}{2em}{1}

Prof A.M.~Muasya       \hfill muthama.muasya@uct.ac.za, (+27) 21 650 3725 \break
Department of Biological Sciences, UCT                                    \break
Keeper, Bolus Herbarium, UCT

Prof G.A.~Verboom        \hfill tony.verboom@uct.ac.za, (+27) 21 650 3398 \break
Head, Department of Biological Sciences, UCT                              \break
Director, Bolus Herbarium, UCT

Prof M.D.~Cramer       \hfill michael.cramer@uct.ac.za, (+27) 21 650 2444 \break
Department of Biological Sciences, UCT

Prof A.G.~West              \hfill adam.west@uct.ac.za, (+27) 21 650 3628 \break
Department of Biological Sciences, UCT

Associate Prof R.~Altwegg \hfill res.altwegg@uct.ac.za, (+27) 21 650 5750 \break
Department of Statistical Sciences, UCT                                   \break
Head, SEEC, UCT

Prof T.A.J.~Hedderson \hfill terry.hedderson@uct.ac.za, (+27) 21 650 4037 \break
Department of Biological Sciences, UCT                                    \break
Bolus Herbarium, UCT

\hfill

\end{hangparas}

\end{document}
